\newpage
\phantomsection
\label{prf:negation-exchange-sup-inf}

\begin{remark*}[Return]
\hyperref[thm:negation-exchange-sup-inf]{Return to Theorem}
\end{remark*}

\begin{theorem*}[Negation Exchanges Supremum and Infimum]
Let $A \subseteq \mathbb{R}$ be a non-empty set. If $A$ is bounded below, then $\sup(-A) = -\inf A$.
\end{theorem*}

\begin{proof}
\textbf{Professional Standard Proof.}~\\
Let $i = \inf A$. By definition, $i \le a$ for all $a \in A$. Multiplying by $-1$ reverses the inequality, giving $-i \ge -a$ for all $-a \in -A$. Thus, $-i$ is an upper bound of $-A$. To show it is the least upper bound, let $u$ be any upper bound of $-A$. Then $u \ge -a$ for all $a \in A$, which implies $-u \le a$. Thus, $-u$ is a lower bound of $A$. Since $i$ is the greatest lower bound, $i \ge -u$, which implies $-i \le u$. Hence $\sup(-A) = -i = -\inf A$.
\end{proof}

\begin{proof}
\textbf{Detailed Learning Proof.}~\\

\textbf{Step 1. Show $-i$ is an upper bound for the reflected set.}
Let $i = \inf A$. For any element $x \in -A$, there exists an element $a \in A$ such that $x = -a$. By the definition of the infimum, $i \le a$. Multiplying both sides by $-1$ reverses the inequality:
\[ -i \ge -a \]
Since $x = -a$, we have $-i \ge x$ for all $x \in -A$, meaning $-i$ is an upper bound of $-A$.

\textbf{Step 2. Establish "Leastness" using the definition of infimum.}
To prove $-i$ is the \textit{least} upper bound, we assume there is some other upper bound $u$ for $-A$. By definition, $u \ge x$ for all $x \in -A$. This means $u \ge -a$ for all $a \in A$.
Multiplying by $-1$ again:
\[ -u \le a \quad \forall a \in A \]
This identifies $-u$ as a lower bound for the original set $A$.

\textbf{Step 3. Conclude by comparing the bounds.}
We know $i$ is the \textbf{Greatest Lower Bound} of $A$. Since $-u$ is just \textit{a} lower bound, it cannot exceed the greatest one:
\[ i \ge -u \]
Multiplying by $-1$ a final time flips the inequality back:
\[ -i \le u \]
Because $-i$ is less than or equal to any arbitrary upper bound $u$, it must be the least upper bound of $-A$.
\end{proof}

\begin{remark*}[Strategy: The Choice of Leastness Method]
In this proof, we establish the "leastness" property by comparing $-i$ against an arbitrary upper bound $u$ (the "Side A" method) rather than using the $\varepsilon$-characterization ("Side B"). While the $\varepsilon$-characterization is mathematically equivalent, it often requires "double negation" logic here—showing that $a < i + \varepsilon$ implies $-a > -i - \varepsilon$. When a proof already involves a set transformation like negation, it is often procedurally cleaner to use the direct definition of the supremum to avoid the mental overhead of tracking flipped inequalities across a tolerance variable.
\end{remark*}

\begin{remark*}[Proof structure]
The proof follows a direct argument by showing $- \inf A$ satisfies both conditions of the supremum for set $-A$: it is an upper bound, and it is less than or equal to any other upper bound.
\end{remark*}

\begin{remark*}[Dependencies]~\\
\begin{itemize}
  \item \hyperref[pred:infimum]{Definition of Infimum}
  \item \hyperref[pred:supremum]{Definition of Supremum}
  \item \hyperref[ax:inequality-negation]{Order Reversal under Negation}
\end{itemize}
\end{remark*}